\section{Gabarito e resoluções comentadas — Competências Digitais}

\begin{solucao}{02.01}\textbf{CERTO.} TLS protege o transporte e autentica o domínio conforme certificado, não a veracidade do conteúdo.\end{solucao}
\begin{solucao}{02.02}\textbf{ERRADO.} O modo privado reduz rastros locais; rede, provedor e sites ainda podem observar metadados e acessos.\end{solucao}
\begin{solucao}{02.03}\textbf{CERTO.} Link público pode ser repassado e não vincula o acesso a identidades específicas.\end{solucao}
\begin{solucao}{02.04}\textbf{CERTO.} A LGPD prevê diversas bases legais; consentimento não é fundamento universal.\end{solucao}
\begin{solucao}{02.05}\textbf{CERTO.} A anonimização efetiva altera o regime, observadas as condições legais e possibilidade razoável de reversão.\end{solucao}
\begin{solucao}{02.06}\textbf{CERTO.} Transparência convive com proteção de dados pessoais e sigilos legais.\end{solucao}
\begin{solucao}{02.07}\textbf{ERRADO.} Neutralidade trata tratamento isonômico nos termos legais, não identidade absoluta de desempenho.\end{solucao}
\begin{solucao}{02.08}\textbf{CERTO.} Reutilização permite credential stuffing após vazamento em outro serviço.\end{solucao}
\begin{solucao}{02.09}\textbf{CERTO.} MFA reduz risco, mas há phishing em tempo real, fadiga de push e outros vetores.\end{solucao}
\begin{solucao}{02.10}\textbf{CERTO.} Histórico melhora rastreabilidade; permissões continuam sendo controle separado.\end{solucao}
\begin{solucao}{02.11}\textbf{B.} Pedido sensível e incomum deve ser validado por canal independente antes de qualquer envio.\end{solucao}
\begin{solucao}{02.12}\textbf{A.} Expressão exata e restrição por domínio reduzem ruído e priorizam fonte oficial.\end{solucao}
\begin{solucao}{02.13}\textbf{CERTO.} Fonte secundária explica; fonte oficial confirma o texto normativo vigente.\end{solucao}
\begin{solucao}{02.14}\textbf{B.} Permissões devem ser proporcionais à finalidade; solicitações excessivas são sinal de risco.\end{solucao}
\begin{solucao}{02.15}\textbf{B.} Compartilhamento nominal e menor privilégio reduzem exposição e aumentam accountability.\end{solucao}
\begin{solucao}{02.16}\textbf{CERTO.} Gravação contém imagem, voz e conteúdo e precisa de governança de finalidade, acesso e retenção.\end{solucao}
\begin{solucao}{02.17}\textbf{B.} O poder público pode tratar dados com base em outras hipóteses legais previstas.\end{solucao}
\begin{solucao}{02.18}\textbf{B.} Controlador toma decisões essenciais; operador atua segundo instruções do controlador.\end{solucao}
\begin{solucao}{02.19}\textbf{CERTO.} O papel jurídico depende das funções efetivamente exercidas.\end{solucao}
\begin{solucao}{02.20}\textbf{B.} Transparência deve observar necessidade e proteção de dados que não agregam à finalidade pública.\end{solucao}
\begin{solucao}{02.21}\textbf{CERTO.} Ativa independe de pedido; passiva atende solicitação.\end{solucao}
\begin{solucao}{02.22}\textbf{A.} Neutralidade envolve isonomia de tratamento com exceções previstas/regulamentadas.\end{solucao}
\begin{solucao}{02.23}\textbf{A.} Texto visual pode mascarar domínio real, técnica comum de phishing.\end{solucao}
\begin{solucao}{02.24}\textbf{CERTO.} O encurtamento dificulta inspeção imediata do destino.\end{solucao}
\begin{solucao}{02.25}\textbf{A.} Coautoria e histórico reduzem versões conflitantes e preservam alterações.\end{solucao}
\begin{solucao}{02.26}\textbf{C.} I e III corretas; HTTPS não comprova factualidade.\end{solucao}
\begin{solucao}{02.27}\textbf{CERTO.} O certificado pode ser válido para o domínio fraudulento.\end{solucao}
\begin{solucao}{02.28}\textbf{B.} Necessidade limita tratamento ao mínimo compatível com a finalidade.\end{solucao}
\begin{solucao}{02.29}\textbf{B.} Coletar atributos sem relação com o serviço viola minimização/necessidade e aumenta risco.\end{solucao}
\begin{solucao}{02.30}\textbf{CERTO.} Biométrico vinculado a pessoa integra categoria sensível.\end{solucao}
\begin{solucao}{02.31}\textbf{B.} Quando juridicamente cabível, acesso pode ser preservado com proteção dos dados não necessários.\end{solucao}
\begin{solucao}{02.32}\textbf{B.} A análise envolve proteção integral, dignidade, privacidade e contexto.\end{solucao}
\begin{solucao}{02.33}\textbf{CERTO.} A proteção no ambiente digital é mais ampla que privacidade informacional.\end{solucao}
\begin{solucao}{02.34}\textbf{A.} Registros não são abertos indistintamente; o acesso segue o regime legal aplicável.\end{solucao}
\begin{solucao}{02.35}\textbf{CERTO.} Decisões formais precisam de preservação em canal/sistema institucional adequado.\end{solucao}
\begin{solucao}{02.36}\textbf{A.} Compartilhar apenas a janela e controlar notificações reduz exposição acidental.\end{solucao}
\begin{solucao}{02.37}\textbf{CERTO.} Extensões podem receber permissões amplas sobre páginas e dados.\end{solucao}
\begin{solucao}{02.38}\textbf{B.} Verificação em fonte oficial e confronto de data/conteúdo reduz desinformação.\end{solucao}
\begin{solucao}{02.39}\textbf{A.} Cópias paralelas quebram fonte única e dificultam saber qual é a versão válida.\end{solucao}
\begin{solucao}{02.40}\textbf{CERTO.} Gerenciador ajuda a criar credenciais únicas, mas a conta/cofre precisa de proteção forte.\end{solucao}
\begin{solucao}{02.41}\textbf{CERTO.} A posse do link é um controle fraco e pode escapar do conjunto esperado de usuários.\end{solucao}
\begin{solucao}{02.42}\textbf{B.} Resposta deve conter, preservar evidência e avaliar escopo/risco conforme processo institucional.\end{solucao}
\begin{solucao}{02.43}\textbf{A.} Links antigos sem owner exigem inventário, expiração e recertificação.\end{solucao}
\begin{solucao}{02.44}\textbf{CERTO.} Migração sem revisão transfere a configuração ruim.\end{solucao}
\begin{solucao}{02.45}\textbf{B.} Transparência da remuneração não justifica dados bancários ou endereço residencial sem finalidade.\end{solucao}
\begin{solucao}{02.46}\textbf{CERTO.} As leis devem ser harmonizadas caso a caso.\end{solucao}
\begin{solucao}{02.47}\textbf{B.} MFA tem qualidade e resistência diferentes; aprovação indevida mostra que desenho/UX precisam ser avaliados.\end{solucao}
\begin{solucao}{02.48}\textbf{B.} O achado trata o controle ausente e o risco, sem atribuir o problema à nuvem em si.\end{solucao}
\begin{solucao}{02.49}\textbf{CERTO.} Competência digital inclui julgamento, verificação e compreensão de limites.\end{solucao}
\begin{solucao}{02.50}\textbf{B.} O perfil institucional representa o órgão; confundir opinião pessoal com posição oficial compromete impessoalidade e credibilidade.\end{solucao}
